\section{Gambaran Umum Program}
Implementasi Milestone 2 melanjutkan struktur program dari Milestone 1 dengan 
menambahkan komponen parser pada pipeline kompilasi. Alur utama program berada 
pada berkas \texttt{src/main.cpp}, yang bertugas membaca masukan, menjalankan lexer 
bila masukan berupa source code, atau langsung memuat token bila masukan sudah berbentuk 
token stream. Setelah token tersedia, parser membangun parse tree dan hasilnya dicetak ke 
terminal serta dapat disimpan ke file keluaran. Desain ini menjaga agar pipeline tetap 
fleksibel untuk kebutuhan pengujian: program dapat diuji langsung dari token stream hasil 
Milestone 1 maupun dari source code Arion.

Modul lexer tetap digunakan untuk mengubah source code menjadi daftar token, sedangkan modul
 \texttt{token\_stream\_reader} berfungsi mengonversi file token dump menjadi objek token 
 yang bisa diproses parser. Modul parser mengandung implementasi Recursive Descent untuk 
 seluruh non-terminal pada grammar Arion, dan modul parse tree menyimpan struktur node 
 beserta mekanisme pencetakan pohon. Struktur ini mencerminkan spesifikasi pada \texttt{spek2.md} 
 dan menjaga pemisahan tanggung jawab antar komponen, sehingga parser dapat fokus pada verifikasi 
 sintaks dan penyusunan parse tree.

\section{Algoritma Parser}
Algoritma parser yang digunakan adalah Recursive Descent dengan satu fungsi untuk setiap 
non-terminal grammar. Parser membaca token secara berurutan dengan strategi top-down, 
dimulai dari simbol awal \texttt{<program>}. Pada setiap fungsi, parser memeriksa token saat 
ini menggunakan lookahead sederhana, memilih produksi yang relevan, lalu mengonsumsi token 
jika cocok. Jika token tidak sesuai, parser menghentikan proses dan mengeluarkan pesan kesalahan 
yang memuat lokasi baris dan kolom.

Logika pencocokan token dibangun melalui tiga operasi utama: \texttt{check} untuk melakukan 
inspeksi token tanpa menggeser posisi, \texttt{consume} untuk memvalidasi dan mengonsumsi token 
yang diharapkan, serta \texttt{advance} untuk berpindah token ketika tipe sudah cocok. Ketika 
terjadi mismatch, parser memanggil \texttt{throwSyntaxError} agar pesan error konsisten di seluruh 
aturan. Pola ini membuat alur parsing mudah diikuti, karena setiap fungsi hanya berfokus pada aturan 
produksinya sendiri.

\section{Struktur Direktori Program}
Struktur direktori program pada Milestone 2 mengikuti pola yang sama dengan Milestone 1, dengan penambahan modul parser. Penjelasan ringkas struktur direktori utama ditunjukkan berikut.

\begin{itemize}
	\item \textbf{\texttt{src/}}. Direktori utama yang berisi seluruh source code program dalam bahasa C++. Direktori ini terbagi menjadi beberapa modul.
		  \begin{itemize}
			  \item \texttt{main.cpp}. Entry point program. Bertanggung jawab membaca input, menentukan apakah input berupa token stream atau source code, menjalankan lexer bila diperlukan, memanggil parser, serta mencetak dan menyimpan parse tree.

			  \item \texttt{lexer/}. Subdirektori lexer dan token, berisi:
					\begin{itemize}
						\item \texttt{token.h} dan \texttt{token.cpp}. Definisi \texttt{TokenType}, struktur \texttt{Token}, serta fungsi utilitas format token.
						\item \texttt{lexer.h} dan \texttt{lexer.cpp}. Implementasi DFA lexer untuk menghasilkan token dari source code.
						\item \texttt{token\_stream\_reader.h} dan \texttt{token\_stream\_reader.cpp}. Parser token stream yang mengubah output Milestone 1 menjadi objek token.
					\end{itemize}

			  \item \texttt{parser/}. Subdirektori parser, berisi:
					\begin{itemize}
						\item \texttt{parser.h} dan \texttt{parser.cpp}. Implementasi Recursive Descent Parser untuk seluruh non-terminal grammar Arion.
						\item \texttt{parse\_tree\_node.h} dan \texttt{parse\_tree\_node.cpp}. Definisi node parse tree dan fungsi pencetakan pohon.
						\item \texttt{parse\_tree\_utils.h} dan \texttt{parse\_tree\_utils.cpp}. Kumpulan utilitas traversal dan pencarian node parse tree.
					\end{itemize}
		  \end{itemize}

	\item \textbf{\texttt{doc/}}. Direktori dokumentasi proyek yang berisi laporan Milestone 2 dalam format LaTeX.

	\item \textbf{\texttt{test/}}. Direktori test case untuk pengujian program. Di dalamnya terdapat subdirektori \texttt{milestone-2/} yang menyimpan input dan output pengujian parser.

	\item \textbf{\texttt{Makefile}}. Berisi rule kompilasi program.

	\item \textbf{\texttt{README.md}}. Dokumentasi ringkas program, cara build/run, serta pembagian tugas.
\end{itemize}

\section{Implementasi Recursive Descent Parser}
Implementasi Recursive Descent diwujudkan dalam \texttt{parser.cpp} dan \texttt{parser.h}. Modul ini membangun parse tree dengan memetakan setiap non-terminal grammar ke fungsi rekursif, serta menegakkan aturan token yang diharapkan melalui mekanisme \texttt{check}, \texttt{consume}, dan \texttt{throwSyntaxError}.

\subsection{Alur Utama Parsing}
Parser dibentuk dari token hasil lexer atau token stream. Konstruktor \texttt{Parser} menyaring token \texttt{COMMENT} dan menolak token \texttt{ERROR}/\texttt{UNKNOWN} agar masalah leksikal tidak merusak tahap parsing. Entry point parsing adalah \texttt{parseProgram()}, yang menyusun urutan \texttt{program-header}, \texttt{declaration-part}, \texttt{compound-statement}, dan token \texttt{period}. Setelah itu parser memastikan tidak ada token tersisa. Skema ini mengikuti aturan produksi utama pada grammar Arion, menjaga determinisme, dan memastikan program selalu ditutup dengan \texttt{period}.

\paragraph{Flowchart}
Diagram berikut merangkum alur kerja parser dari input hingga parse tree terbentuk.

\begin{figure}[H]
\centering
\resizebox{\textwidth}{!}{%
\begin{tikzpicture}[
    node distance=1.3cm,
    >=stealth,
    startend/.style={ellipse, draw, minimum width=3.2cm, minimum height=0.8cm, align=center},
    process/.style={rectangle, draw, rounded corners, minimum width=3.6cm, minimum height=0.8cm, align=center},
    decision/.style={diamond, draw, aspect=2.2, minimum width=3.2cm, align=center},
    connector/.style={circle, draw, minimum size=0.55cm, inner sep=0pt}
]
\node (input) [startend] {Input text};
\node (stream) [decision, below=of input] {Token stream?};
\node (lexer) [process, left=3.6cm of stream] {Lexer::tokenize};
\node (tokenstream) [process, right=3.6cm of stream] {TokenStreamReader::tryRead};
\node (parser) [process, below=2.1cm of stream] {Parser constructor};
\node (parseprogram) [process, below=of parser] {parseProgram};
\node (header) [process, below=of parseprogram] {parseProgramHeader};
\node (decl) [process, below=of header] {parseDeclarationPart};
\node (compound) [process, below=of decl] {parseCompoundStatement};
\node (period) [process, below=of compound] {consume period};
\node (endcheck) [decision, below=of period] {isAtEnd?};
\node (ok) [startend, below=of endcheck] {Return parse tree};
\node (error) [process, right=3.6cm of endcheck] {throwSyntaxError};

\draw[->] (input) -- (stream);
\draw[->] (stream.west) -- node[above] {no} (lexer.east);
\draw[->] (stream.east) -- node[above] {yes} (tokenstream.west);
\draw[->] (lexer) |- (parser);
\draw[->] (tokenstream) |- (parser);
\draw[->] (parser) -- (parseprogram);
\draw[->] (parseprogram) -- (header);
\draw[->] (header) -- (decl);
\draw[->] (decl) -- (compound);
\draw[->] (compound) -- (period);
\draw[->] (period) -- (endcheck);
\draw[->] (endcheck) -- node[left] {yes} (ok);
\draw[->] (endcheck.east) -- node[above] {no} (error.west);
\end{tikzpicture}%
}
\caption{Flowchart alur utama parser}
\label{fig:parser-flowchart}
\end{figure}

\paragraph{Pseudocode}
Berikut ringkasan algoritma dari \texttt{parseProgram()} berdasarkan implementasi aktual.

\begin{tcolorbox}[title=Pseudocode parseProgram,colback=gray!5,colframe=black!60,boxrule=0.5pt,arc=2mm,breakable]
\begin{verbatim}
function parseProgram(tokens):
    parser = Parser(tokens)
    node = ProgramNode()
    node.add(parseProgramHeader())
    node.add(parseDeclarationPart())
    node.add(parseCompoundStatement())
    node.add(terminal(consume(PERIOD)))
    if not isAtEnd():
        throwSyntaxError("end of input")
    return node
\end{verbatim}
\end{tcolorbox}

\subsection{Pemilihan Produksi Dengan Lookahead}
Pemilihan produksi dilakukan dengan lookahead satu token melalui \texttt{check} dan variasi \texttt{check(..., offset)}. Strategi ini dipakai ketika satu token awal bisa mengarah ke beberapa produksi. Contohnya, pada \texttt{parseStatement} parser membedakan assignment dan procedure/function call hanya dari token \texttt{ident} yang diikuti \texttt{becomes}, \texttt{lbrack}, atau \texttt{period}. Pola yang sama dipakai pada \texttt{parseFactor} untuk membedakan pemanggilan prosedur/fungsi, akses variabel berkomponen, dan ident biasa. Dengan lookahead satu token, parser tidak memerlukan backtracking sehingga tetap sederhana dan efisien.

\begin{tcolorbox}[title=Pseudocode lookahead,colback=gray!5,colframe=black!60,boxrule=0.5pt,arc=2mm,breakable]
\begin{verbatim}
function check(type, offset = 0):
    token = peek(offset)
    return token != null and token.type == type

function parseStatement():
    if check(IDENT):
        if check(BECOMES, 1) or check(LBRACK, 1) or check(PERIOD, 1):
            return parseAssignmentStatement()
        else:
            return parseProcedureFunctionCall()
    ...
\end{verbatim}
\end{tcolorbox}

\subsection{Struktur Fungsi Per Non-Terminal}
Non-terminal deklarasi seperti 
\texttt{const-declaration}, 
\texttt{type-declaration}, dan \texttt{var-declaration} diimplementasikan 
sebagai loop agar dapat menerima jumlah deklarasi yang bervariasi. 
Non-terminal statement seperti \texttt{if}, \texttt{case}, \texttt{while}, 
\texttt{repeat}, dan \texttt{for} dipecah ke fungsi khusus agar struktur 
grammar tercermin langsung di dalam kode. Dengan pola ini, urutan pemanggilan 
fungsi parser menyerupai proses derivasi grammar secara eksplisit.

Fungsi dispatcher seperti \texttt{parseType()} dan \texttt{parseStatement()} memilih cabang produksi berdasarkan token saat ini. Non-terminal yang opsional (misalnya statement kosong) ditangani dengan pengecekan follow set sederhana, sehingga parser tetap konsisten dengan grammar tanpa memaksakan token yang tidak ada.

\begin{tcolorbox}[title=Pseudocode struktur non-terminal,colback=gray!5,colframe=black!60,boxrule=0.5pt,arc=2mm,breakable]
\begin{verbatim}
function parseDeclarationPart():
    node = DeclarationPartNode()
    while check(CONSTSY):
        node.add(parseConstDeclaration())
    while check(TYPESY):
        node.add(parseTypeDeclaration())
    while check(VARSY):
        node.add(parseVarDeclaration())
    while check(PROCEDURESY) or check(FUNCTIONSY):
        node.add(parseSubprogramDeclaration())
    return node

function parseStatement():
    node = StatementNode()
    if check(IDENT):
        if check(BECOMES, 1) or check(LBRACK, 1) or check(PERIOD, 1):
            node.add(parseAssignmentStatement())
        else:
            node.add(parseProcedureFunctionCall())
        return node
    if check(BEGINSY): return node.add(parseCompoundStatement())
    if check(IFSY): return node.add(parseIfStatement())
    if check(CASESY): return node.add(parseCaseStatement())
    if check(WHILESY): return node.add(parseWhileStatement())
    if check(REPEATSY): return node.add(parseRepeatStatement())
    if check(FORSY): return node.add(parseForStatement())
    if isStatementFollow(): return node.add(EmptyStatementNode())
    throwSyntaxError("statement")
\end{verbatim}
\end{tcolorbox}

\subsection{Hierarki Ekspresi dan Prioritas Operator}
Ekspresi dipecah menjadi empat level: 
\texttt{expression}, 
\texttt{simple-expression}, \texttt{term}, dan \texttt{factor}. 
Pemisahan ini memastikan prioritas operator diproses dengan benar. 
Operator relasional ditangani pada \texttt{expression}, operator 
aditif pada \texttt{simple-expression}, dan operator multiplikatif 
pada \texttt{term}. \texttt{factor} menangani literal, variabel, pemanggilan 
prosedur/fungsi, serta ekspresi di dalam tanda kurung.

Struktur ini menjaga agar ekspresi seperti \texttt{a + b * c} diproses sesuai precedence, karena \texttt{parseTerm()} selalu menyelesaikan operasi multiplikatif sebelum \texttt{parseSimpleExpression()} menggabungkannya dengan operator aditif.

\begin{tcolorbox}[title=Pseudocode hierarki ekspresi,colback=gray!5,colframe=black!60,boxrule=0.5pt,arc=2mm,breakable]
\begin{verbatim}
function parseExpression():
    node = ExpressionNode()
    node.add(parseSimpleExpression())
    if next is relop:
        node.add(parseRelationalOperator())
        node.add(parseSimpleExpression())
    return node

function parseSimpleExpression():
    node = SimpleExpressionNode()
    if next is PLUS or MINUS:
        node.add(terminal(advance()))
    node.add(parseTerm())
    while next is PLUS or MINUS or ORSY:
        node.add(parseAdditiveOperator())
        node.add(parseTerm())
    return node

function parseTerm():
    node = TermNode()
    node.add(parseFactor())
    while next is TIMES or RDIV or IDIV or IMOD or ANDSY:
        node.add(parseMultiplicativeOperator())
        node.add(parseFactor())
    return node

function parseFactor():
    if next is IDENT:
        if lookahead is LPARENT:
            return parseProcedureFunctionCall()
        if lookahead is LBRACK or PERIOD:
            return parseVariable()
        return terminal(advance())
    if next is literal:
        return terminal(advance())
    if next is LPARENT:
        advance(); node = parseExpression(); consume(RPARENT)
		return node
    if next is NOTSY:
        advance(); return parseFactor()
    throwSyntaxError("factor")
\end{verbatim}
\end{tcolorbox}

\subsection{Variabel Dengan Komponen Berantai}
Parser mendukung akses array dan record melalui \texttt{parseVariableWithBase} yang membungkus variabel dasar ke dalam \texttt{component-variable} secara rekursif. Variabel dasar selalu diawali \texttt{ident}. Setiap suffix berupa index \texttt{[ ... ]} atau akses field \texttt{.ident} menghasilkan node baru yang menampung base variable sebelumnya, sehingga konstruk seperti \texttt{arr[1].field} diuraikan sebagai rantai node bertingkat. Untuk Milestone 2, isi \texttt{index-list} dibatasi pada \texttt{intcon}, \texttt{charcon}, atau \texttt{ident}.

\begin{tcolorbox}[title=Pseudocode variable berantai,colback=gray!5,colframe=black!60,boxrule=0.5pt,arc=2mm,breakable]
\begin{verbatim}
function parseVariable():
    base = VariableNode()
    base.add(terminal(consume(IDENT)))
    return parseVariableWithBase(base)

function parseVariableWithBase(base):
    if not (next is LBRACK or next is PERIOD):
        return base
    wrapped = VariableNode()
    wrapped.add(parseComponentVariable(base))
    return parseVariableWithBase(wrapped)

function parseComponentVariable(base):
    node = ComponentVariableNode()
    node.add(base)
    if next is LBRACK:
        node.add(terminal(advance()))
        node.add(parseIndexList())
        node.add(terminal(consume(RBRACK)))
        return node
    node.add(terminal(consume(PERIOD)))
    node.add(terminal(consume(IDENT)))
    return node

function parseIndexList():
    node = IndexListNode()
    require next in {INTCON, CHARCON, IDENT}
    node.add(terminal(advance()))
    while next is COMMA:
        node.add(terminal(advance()))
        require next in {INTCON, CHARCON, IDENT}
        node.add(terminal(advance()))
    return node
\end{verbatim}
\end{tcolorbox}

\subsection{Penanganan Error dan Validasi Token}
Penanganan error dipusatkan pada \texttt{throwSyntaxError}, yang menghasilkan pesan berformat konsisten dengan informasi baris dan kolom. Parser juga memfilter token komentar serta menolak token \texttt{ERROR} dan \texttt{UNKNOWN} sejak konstruksi parser, sehingga masalah leksikal tidak merusak tahap parsing. Dengan cara ini, parser berhenti secara deterministik ketika menemukan pelanggaran grammar dan memberikan diagnosis yang jelas.

\section{Implementasi Kode}
Bagian ini menjelaskan isi setiap file dengan memisahkan penjelasan per snippet penting agar struktur implementasi mudah ditelusuri.

\subsection{\texttt{main.cpp}}
\subsubsection{Alur Utama Eksekusi}
Snippet berikut menunjukkan alur utama: validasi argumen, pembacaan input, pemilihan jalur token stream atau lexer, pemanggilan parser, dan pencetakan output.

\begin{lstlisting}[language=C++, numbers=left, numberstyle=\tiny, breaklines=true, breakatwhitespace=true, caption={Alur utama program pada main.cpp}, label={lst:main-flow}]
if (argc < 2) {
	std::cerr << "Usage: " << argv[0] << " <input.txt> [output.txt]" << std::endl;
	return 1;
}

Lexer lexer(source);
	if (!TokenStreamReader::tryRead(source, tokens, tokenStreamError)) {
		std::cerr << "Error: " << tokenStreamError << std::endl;
		return 1;
	}
} else {
	tokens = lexer.tokenize();
	if (lexer.hasErrors()) {
		return 1;
	}
}

Parser parser(tokens);
std::unique_ptr<ProgramNode> parseTree = parser.parseProgram();
parseTree->print(outputBuffer);
\end{lstlisting}

\subsection{\texttt{lexer/token.h}}
\subsubsection{Enum \texttt{TokenType}}
Enum ini mendefinisikan seluruh jenis token yang dipakai lexer dan parser.

\begin{lstlisting}[language=C++, numbers=left, numberstyle=\tiny, breaklines=true, breakatwhitespace=true, caption={Enum TokenType pada token.h}, label={lst:token-enum-m2}]
enum class TokenType {
	INTCON, REALCON, CHARCON, STRING,
	NOTSY, ANDSY, ORSY,
	PLUS, MINUS, TIMES, IDIV, RDIV, IMOD,
	EQL, NEQ, GTR, GEQ, LSS, LEQ,
	LPARENT, RPARENT, LBRACK, RBRACK, COMMA,
	SEMICOLON, PERIOD, COLON, BECOMES,
	CONSTSY, TYPESY, VARSY, FUNCTIONSY, PROCEDURESY,
	ARRAYSY, RECORDSY, PROGRAMSY,
	IDENT, BEGINSY, IFSY, CASESY, REPEATSY,
	WHILESY, FORSY, ENDSY, ELSESY, UNTILSY,
	OFSY, DOSY, TOSY, DOWNTOSY, THENSY,
	COMMENT, UNKNOWN, ERROR
};
\end{lstlisting}

\subsubsection{Struct \texttt{Token}}
Struct ini menyimpan informasi token beserta lokasi baris dan kolom.

\begin{lstlisting}[language=C++, numbers=left, numberstyle=\tiny, breaklines=true, breakatwhitespace=true, caption={Struct Token pada token.h}, label={lst:token-struct-m2}]
struct Token {
	TokenType type;
	std::string value;
	int line;
	int column;
};
\end{lstlisting}

\subsection{\texttt{lexer/token.cpp}}
\subsubsection{Format Output Token}
Fungsi berikut memastikan token dicetak sesuai format spesifikasi.

\begin{lstlisting}[language=C++, numbers=left, numberstyle=\tiny, breaklines=true, breakatwhitespace=true, caption={Implementasi formatToken pada token.cpp}, label={lst:format-token-m2}]
std::string formatToken(const Token& token) {
	if (tokenHasValue(token.type)) {
		return tokenTypeToString(token.type) + 
            " (" + token.value + ")";
	}
	return tokenTypeToString(token.type);
}
\end{lstlisting}

\subsection{\texttt{lexer/token\_stream\_reader.h}}
\subsubsection{Antarmuka Pembaca Token Stream}
Header ini mengekspos fungsi utilitas untuk mendeteksi dan membaca token stream.

\begin{lstlisting}[language=C++, numbers=left, numberstyle=\tiny, breaklines=true, breakatwhitespace=true, caption={Deklarasi TokenStreamReader}, label={lst:token-stream-api}]
namespace TokenStreamReader {
	bool looksLikeTokenStream(const std::string &text);
	bool tryRead(const std::string &text, 
        std::vector<Token> &tokens,
        std::string &errorMessage);
}
\end{lstlisting}

\subsection{\texttt{lexer/token\_stream\_reader.cpp}}
\subsubsection{Pemetaan String ke TokenType}
Token stream menggunakan label string yang dipetakan ke enum \texttt{TokenType}.

\begin{lstlisting}[language=C++, numbers=left, numberstyle=\tiny, breaklines=true, breakatwhitespace=true, caption={Pemetaan token pada token\_stream\_reader.cpp}, label={lst:token-stream-map}]
const std::unordered_map<std::string, TokenType> &tokenTypeMap()
{
	static const std::unordered_map<std::string, TokenType> mapping = {
		{"intcon", TokenType::INTCON},
		{"realcon", TokenType::REALCON},
		{"charcon", TokenType::CHARCON},
		{"string", TokenType::STRING},
		{"programsy", TokenType::PROGRAMSY},
		{"ident", TokenType::IDENT},
		{"semicolon", TokenType::SEMICOLON},
		{"period", TokenType::PERIOD},
		{"error", TokenType::ERROR},
	};
	return mapping;
}
\end{lstlisting}

\subsubsection{Parsing Token Stream}
Fungsi \texttt{parseSequence} melakukan parsing token stream secara rekursif hingga akhir input.

\begin{lstlisting}[language=C++, numbers=left, numberstyle=\tiny, breaklines=true, breakatwhitespace=true, caption={Bagian parseSequence pada token\_stream\_reader.cpp}, label={lst:token-stream-parse}]
bool parseSequence(std::size_t position, std::vector<Token> &tokens) const
{
	position = skipSeparators(position);
	if (position >= text.size()) {
		tokens.clear();
		return true;
	}
	return parseSequence(nextPosition, tokens);
}
\end{lstlisting}

\subsection{\texttt{parser/parser.h}}
\subsubsection{Kelas \texttt{ParserError}}
Kelas ini menyimpan pesan error dan lokasi sumber error.

\begin{lstlisting}[language=C++, numbers=left, numberstyle=\tiny, breaklines=true, breakatwhitespace=true, caption={ParserError pada parser.h}, label={lst:parser-error}]
class ParserError : public std::runtime_error {
public:
	ParserError(std::string message, int line, int column);
	int line() const;
	int column() const;
private:
	int errorLine;
	int errorColumn;
};
\end{lstlisting}

\subsubsection{Kelas \texttt{Parser}}
Kelas ini mendeklarasikan seluruh fungsi non-terminal dan utilitas parsing.

\begin{lstlisting}[language=C++, numbers=left, numberstyle=\tiny, breaklines=true, breakatwhitespace=true, caption={Deklarasi inti Parser pada parser.h}, label={lst:parser-class}]
class Parser {
public:
	explicit Parser(const std::vector<Token> &tokens);
	std::unique_ptr<ProgramNode> parseProgram();
	std::unique_ptr<ProgramHeaderNode> parseProgramHeader();
private:
	std::vector<Token> parserTokens;
	std::size_t currentIndex;
	bool check(TokenType type) const;
	Token consume(TokenType type, const std::string &expectedName);
	[[noreturn]] void throwSyntaxError(const std::string &expectedName) const;
	// ... deklarasi fungsi non-terminal lainnya ...
};
\end{lstlisting}

\subsection{\texttt{parser/parser.cpp}}
\subsubsection{Entry Point Parsing}
Fungsi berikut memulai parsing dan memastikan program diakhiri token \texttt{period}.

\begin{lstlisting}[language=C++, numbers=left, numberstyle=\tiny, breaklines=true, breakatwhitespace=true, caption={parseProgram pada parser.cpp}, label={lst:parse-program}]
std::unique_ptr<ProgramNode> Parser::parseProgram()
{
	auto programNode = std::make_unique<ProgramNode>();
	programNode->addChild(parseProgramHeader());
	programNode->addChild(parseDeclarationPart());
	programNode->addChild(parseCompoundStatement());
	programNode->addChild(makeTerminalNode(consume(
        TokenType::PERIOD, "period")));

	if (!isAtEnd()) {
		throwSyntaxError("end of input");
	}
	return programNode;
}
\end{lstlisting}

\subsubsection{Pelaporan Syntax Error}
Penanganan error terpusat pada \texttt{throwSyntaxError} agar format pesan konsisten.

\begin{lstlisting}[language=C++, numbers=left, numberstyle=\tiny, breaklines=true, breakatwhitespace=true, caption={throwSyntaxError pada parser.cpp}, label={lst:throw-syntax-error}]
void Parser::throwSyntaxError(const std::string &expectedName) const
{
	std::ostringstream message;
	if (isAtEnd()) {
		int line = parserTokens.empty() ? 1 : parserTokens.back().line;
		int column = parserTokens.empty() ? 1 : parserTokens.back().column;
		message << "Syntax error at line " << line
				<< ", column " << column
				<< ": unexpected end of input, expected " << expectedName;
		throw ParserError(message.str(), line, column);
	}

	const Token &token = peek();
	message << "Syntax error at line " << token.line
			<< ", column " << token.column
			<< ": unexpected token " << tokenTypeToString(token.type)
			<< ", expected " << expectedName;
	throw ParserError(message.str(), token.line, token.column);
}
\end{lstlisting}

\subsubsection{Pemilihan Statement}
Parser membedakan assignment dan procedure/function call berdasarkan lookahead.

\begin{lstlisting}[language=C++, numbers=left, numberstyle=\tiny, breaklines=true, breakatwhitespace=true, caption={Bagian parseStatement pada parser.cpp}, label={lst:parse-statement}]
if (check(TokenType::IDENT)) {
	if (check(TokenType::BECOMES, 1) || check(TokenType::LBRACK, 1) || check(TokenType::PERIOD, 1)) {
		statementNode->addChild(parseAssignmentStatement());
	} else {
		statementNode->addChild(parseProcedureFunctionCall());
	}
	return statementNode;
}
\end{lstlisting}

\subsection{\texttt{parser/parse\_tree\_node.h}}
\subsubsection{Struktur Dasar Parse Tree}
Kelas berikut adalah fondasi semua node pada parse tree.

\begin{lstlisting}[language=C++, numbers=left, numberstyle=\tiny, breaklines=true, breakatwhitespace=true, caption={ParseTreeNode pada parse\_tree\_node.h}, label={lst:parse-tree-node}]
class ParseTreeNode {
public:
	explicit ParseTreeNode(std::string name);
	const std::string &name() const;
	void addChild(std::unique_ptr<ParseTreeNode> child);
	void print(std::ostream &out, const std::string &prefix = "", bool isLast = true) const;
protected:
	virtual std::string renderLabel() const;
private:
	std::string nodeName;
	ParseTreeNode *parentNode;
	std::vector<std::unique_ptr<ParseTreeNode>> childNodes;
};
\end{lstlisting}

\subsubsection{Terminal Node}
Node terminal menyimpan token dan mencetak label token saat output.

\begin{lstlisting}[language=C++, numbers=left, numberstyle=\tiny, breaklines=true, breakatwhitespace=true, caption={ParseTreeTerminalNode pada parse\_tree\_node.h}, label={lst:parse-tree-terminal}]
class ParseTreeTerminalNode : public ParseTreeNode {
public:
	explicit ParseTreeTerminalNode(const Token &token);
	const Token &token() const;
protected:
	std::string renderLabel() const override;
private:
	Token terminalToken;
};
\end{lstlisting}

\subsection{\texttt{parser/parse\_tree\_node.cpp}}
\subsubsection{Pencetakan Tree Dengan Indentasi}
Fungsi \texttt{print} membangun struktur pohon dengan simbol garis.

\begin{lstlisting}[language=C++, numbers=left, numberstyle=\tiny, breaklines=true, breakatwhitespace=true, caption={ParseTreeNode::print pada parse\_tree\_node.cpp}, label={lst:parse-tree-print}]
void ParseTreeNode::print(std::ostream &out, const std::string &prefix, bool isLast) const
{
	out << prefix;
	if (parentNode != nullptr) {
		out << (isLast ? "`-- " : "|-- ");
	}
	out << renderLabel() << '\n';

	const std::string childPrefix = prefix + (parentNode == nullptr ? "" : (isLast ? "    " : "|   "));
	for (std::size_t index = 0; index < childNodes.size(); ++index) {
		const bool childIsLast = index + 1 == childNodes.size();
		childNodes[index]->print(out, childPrefix, childIsLast);
	}
}
\end{lstlisting}

\subsection{\texttt{parser/parse\_tree\_utils.h}}
\subsubsection{Deklarasi Utilitas Traversal}
Header ini menyediakan fungsi bantu untuk inspeksi pohon.

\begin{lstlisting}[language=C++, numbers=left, numberstyle=\tiny, breaklines=true, breakatwhitespace=true, caption={Deklarasi utilitas parse tree}, label={lst:parse-tree-utils-decl}]
namespace ParseTreeUtils {
	bool isRoot(const ParseTreeNode *node);
	bool isLeaf(const ParseTreeNode *node);
	std::vector<ParseTreeNode *> collectAllNodes(ParseTreeNode *node);
	std::vector<ParseTreeNode *> collectTerminalNodes(ParseTreeNode *node);
	std::vector<ParseTreeNode *> collectNodesByName(ParseTreeNode *node, const std::string &name);
}
\end{lstlisting}

\subsection{\texttt{parser/parse\_tree\_utils.cpp}}
\subsubsection{Contoh Traversal Preorder}
Contoh berikut memperlihatkan traversal preorder untuk mengumpulkan seluruh node.

\begin{lstlisting}[language=C++, numbers=left, numberstyle=\tiny, breaklines=true, breakatwhitespace=true, caption={collectAllNodes pada parse\_tree\_utils.cpp}, label={lst:parse-tree-utils-collect}]
std::vector<ParseTreeNode *> collectAllNodes(ParseTreeNode *node)
{
	std::vector<ParseTreeNode *> result;
	if (!node)
		return result;

	result.push_back(node);
	for (auto &child : node->children()) {
		auto childNodes = collectAllNodes(child.get());
		result.insert(result.end(), childNodes.begin(), childNodes.end());
	}
	return result;
}
\end{lstlisting}
